% sec-04: the operation and the advisory boundary (surgical lead).
\section{The Operation and the Advisory Boundary}

The intervention is a staged eight-arm robotic pancreaticoduodenectomy with
perioperative daraxonrasib and an on-premises language model whose output is
advisory only \cite{onpremwhippl,phase1ind}. The model has no path to the arms:
no generated token reaches a motion controller and no software in the advisory
path holds write access to the command channel. Arm-level stop is specified at
3 milliseconds and system-wide stop at 500 milliseconds, measured on the bench
before the first patient and again at every cohort boundary.

\begin{apptable}
\begin{tabularx}{\textwidth}{>{\raggedright\arraybackslash}p{3.4cm} >{\raggedright\arraybackslash}p{3.1cm} >{\raggedright\arraybackslash}X}
\headrow \thead{Parameter} & \thead{Value} & \thead{Source}\\
\midrule
Population & Resectable or borderline-resectable KRAS G12-mutated PDAC & Phase 1 protocol \\
Design and size & Open-label, single-arm, 3+3 escalation, up to 18 treated participants & Phase 1 protocol and IND \\
Co-primary endpoints & 30-day device- or procedure-related SAEs; DLT and MTD or RP2D; tasks completed without unsafe conversion & Phase 1 protocol \\
Follow-up & 28-day screening; 30-day and 90-day safety; 24-month overall survival & Phase 1 protocol \\
\end{tabularx}
\end{apptable}
\tabcap{Trial parameters carried unchanged from the published Phase 1 protocol
and the Investigational New Drug application, so the concept is evaluated
against a protocol that already exists.}
